\section{Zadanie 10}
Niech \(X\) będzie liczbą naturalną wylosowaną w sposób jednostajny z przedziału \([1, n]\),
a \(Y\) liczbą całkowitą wylosowaną w sposób jednostajny z przedziału \([-k, k]\).
Wyznaczyć \(\variance{X}\) oraz \(\variance{Y}\).

\subsection*{Rozwiązanie}
\begin{enumerate}

	\item \textbf{Wariancja zmiennej} \(X\), gdzie \(X\) jest wylosowane jednostajnie z przedziału \([1, n]\):

	      \[
		      \begin{aligned}
			      \ev{X}^2
			       & = \left( \frac{1 + 2 + \dots + n}{n} \right)^2
			      = \left( \frac{n + 1}{2} \right)^2,                               \\[6pt]
			      \ev{X^2}
			       & = \frac{1^2 + 2^2 + \dots + n^2}{n}
			      = \frac{(n + 1)(2n + 1)}{6},                                      \\[6pt]
			      \variance{X}
			       & = \ev{X^2} - \ev{X}^2                                          \\[3pt]
			       & = \frac{(n + 1)(2n + 1)}{6} - \left( \frac{n + 1}{2} \right)^2 \\[3pt]
			       & = \boxed{\frac{n^2 - 1}{12}}.
		      \end{aligned}
	      \]

	\item \textbf{Wariancja zmiennej} \(Y\), gdzie \(Y\) jest wylosowane jednostajnie z przedziału \([-k, k]\):

	      \[
		      \begin{aligned}
			      \ev{Y}
			       & = \frac{-k + (-k+1) + \dots + (k-1) + k}{2k+1}
			      = 0,                                              \\[6pt]
			      \ev{Y^2}
			       & = \frac{(-k)^2 + \dots + k^2}{2k+1}
			      = \frac{1}{2k+1}\sum_{n=-k}^{k} n^2               \\[3pt]
			       & = \frac{2}{2k+1}\sum_{n=1}^{k} n^2
			      = \frac{2}{2k+1} \cdot \frac{k(k+1)(2k+1)}{6}
			      = \frac{k(k+1)}{3},                               \\[6pt]
			      \variance{Y}
			       & = \ev{Y^2} - \ev{Y}^2
			      = \boxed{\frac{k(k+1)}{3}}.
		      \end{aligned}
	      \]

\end{enumerate}
